\clearpage
\SourceLeaf{276}%
% A544-P0276-D01
\renewcommand{\BookRunningTitle}{Démonstration de quelques théorèmes}
\pdfbookmark[0]{Démonstration de quelques théorèmes concernants les périodes}{periodes-classes}
\begin{center}
\Large\bfseries DÉMONSTRATION DE QUELQUES THÉORÈMES CONCERNANTS\par\medskip
\large LES PÉRIODES DES CLASSES DES FORMES BINAIRES DU SECOND DEGRÉ.
\end{center}
\BookRule{.18\textwidth}

% A544-P0276-D02
\pdfbookmark[1]{Théorème I}{periodes-theoreme-I}
\textsc{Théorème. I.}\quad\emph{Le nombre des classes (pr. pr.) d’un même déterminant, qui élevées à la dignité $P^{\mathrm{me}}$, $P$ étant ou un nombre premier ou la puissance d’un nombre premier $=p^\pi$, produisent la classe principale $K$, est égal ou à $1$ ou à une puissance de ce même nombre premier $p$.}

% A544-P0276-D03
\emph{Démonstration.}\quad Soit $(\mathfrak Q)$ le groupe entier de toutes les classes en question et $n$ leur nombre. Puisque la classe principale $K$ est nécessairement contenue dans $(\mathfrak Q)$, le théorème est évident, si elle y est la seule. Mais s’il y en a d’autres, le nombre des classes contenues dans la période de chacune sera une puissance de $p$; soit une d’elles $A$, et supposons que sa période $(\mathfrak A)$ contienne $p^\alpha$ classes, qui seront toutes comprises dans $(\mathfrak Q)$. Or si les classes de cette période $(\mathfrak A)$ épuisent $(\mathfrak Q)$, on aura $p^\alpha=n$, et le théorème sera démontré; sinon, soit $B$ une classe quelconque de $(\mathfrak Q)$ non contenue dans $(\mathfrak A)$, et supposons que sa période soit développée jusqu’à ce qu’on y parvienne à une classe $bB$, qui soit en même temps parmi les classes de $(\mathfrak A)$, ce qui doit nécessairement arriver, parceque du moins la classe principale est commune à cette période et à $(\mathfrak A)$. Or supposant que $bB$ soit la \emph{première} classe dans la période de $B$ commune à $(\mathfrak A)$, ou $b$ le plus petit possible, je dis

% A544-P0276-D04
$1^0$.\quad Que $b$ sera une puissance de $p$. Car il est évident qu’en faisant $b=p^\beta h$, $bB=iA$ et $hk\equiv1\pmod{p^\pi}$ (ce qui se pourra) on aura $kbB=p^\beta hkB=p^\beta B=ikA$,
\SourceLeaf{277}%
% A544-P0277-D01
c’est à dire que $p^\beta B$ sera aussi parmi les classes de $(\mathfrak A)$, d’où il s’ensuit que $h=1$ et $b=p^\beta$.

% A544-P0277-D02
$2^0$.\quad Qu’en désignant les classes $K,B,2B\dotsfour(b-1)B$ par $(\mathfrak B)$, toutes les compositions d’une classe de $(\mathfrak A)$ avec une classe de $(\mathfrak B)$ donneront $p^{\alpha+\beta}$ classes différentes. Car en supposant $mA+nB=m'A+n'B$ et $n=n'$, on aura nécessairement $m=m'$; si $n>n'$, on aura $(n-n')B=(m'-m)A$, ce qui est impossible, si l’on n’a pas $n=n'$.

% A544-P0277-D03
$3^0$.\quad Que ces $p^{\alpha+\beta}$ classes différentes seront comprises sous $(\mathfrak Q)$, ce qui est évident.

% A544-P0277-D04
Or, si ces $p^{\alpha+\beta}$ classes épuisent $(\mathfrak Q)$, le théorème est démontré; sinon, on choisira nne autre classe de $(\mathfrak Q)$ non contenue parmi celles-là, savoir $C$; on continuera sa période jusqu’à ce qu’on y parvienne à une classe déjà comprise sous les classes composées de $(\mathfrak A)$ et $(\mathfrak B)$. Par un raisonnement semblable au précédent\textsuperscript{'} on démontrera, que l’exposant de cette classe doit être une puissance de $p$, $=p^\gamma$, et que la composition des $p^\gamma$ classes premières de la période de $C$ avec les $p^{\alpha+\beta}$ classes déjà trouvées donnera $p^{\alpha+\beta+\gamma}$ classes différentes toutes comprises dans $(\mathfrak Q)$. Si ces classes n’épuisent pas encore $(\mathfrak Q)$, on traitera de la même manière une quatrième classe $D$ etc., et il est évident que $(\mathfrak Q)$ étant formé d’un nombre fini de classes, ces opérations finiront aussi et qu’on aura $n$ égal à une puissance de $p$.\quad C. Q. F. D.

% A544-P0277-D05
\par\bigskip
\pdfbookmark[1]{Théorème II}{periodes-theoreme-II}
\textsc{Théorème. II.}\quad\emph{Le nombre de toutes les classes du genre principal étant exprimé par $a^\alpha b^\beta c^\gamma$ etc., $a,b,c$, dénotant des nombres premiers différents, il y aura dans ce genre $a^\alpha,b^\beta,c^\gamma$ etc. classes, qui étant élevées à la dignité $a^\alpha,b^\beta,c^\gamma$ etc. resp. produisent la classe principale.}

% A544-P0277-D06
\emph{Démonstration}\quad Soient $A,A',A''$ etc. toutes les classes qui élevées à la dignité $a^\alpha$ produisent $K$ et $(\mathfrak A)$ leur totalité; de même $B,B',B''$ etc. $(\mathfrak B)$, $C,C',C''$, $(\mathfrak C)$ etc. etc. Je dis que de la composition de toutes les classes de $(\mathfrak A)$ avec toutes les classes de $(\mathfrak B)$ avec toutes les classes de $(\mathfrak C)$ etc. il proviendra des classes différentes entre elles. Car si $A+B+C\dots=A'+B'+C'\dots$ etc., on aura, en faisant $A-A'=A''$, $B-B'=B''$ etc.,
\[
A''+B''+C''\text{ etc.}=K
\]
donc élevant à la dignité $b^\beta c^\gamma$ etc., $(b^\beta c^\gamma\dots)A''=K$, d’où il s’ensuit facilement
\SourceLeaf{278}%
% A544-P0278-D01
$A''=K$ et $A=A'$ et de la même manière on aura $B=B'$, $C=C'$ etc. Soit la totalité de ces classes $=(S)$. De plus il est clair que toutes ces classes seront du genre principal. Enfin il ne peut exister aucune classe dans le genre principal qui ne soit comprise sous $(S)$. Soit\ .\ .\ .
\BookRule{.34\textwidth}

% A544-P0278-D02
\par\bigskip
\pdfbookmark[1]{Bemerkung}{periodes-bemerkung}
\begin{center}\textsc{Bemerkung.}\end{center}
\begingroup\small
Dieses im Jahre 1801 geschriebene Fragment bezieht sich auf Disq. Arithm. art. 306, \textsc{ix}. Das Wort \emph{dignité} wird hier in einem sonst nicht üblichen Sinne gebraucht.
\begin{flushright}\textsc{Stern}.\end{flushright}
\endgroup
\BookRule{.34\textwidth}

\clearpage
\SourceLeaf{279}%
% A544-P0279-D01
\renewcommand{\BookRunningTitle}{De nexu inter multitudinem classium}
\renewcommand{\Article}[1]{\par\medskip\Needspace{4\baselineskip}\pdfbookmark[1]{Art. #1}{nexus-prior-art-#1}\centerline{\textbf{#1.}}\nobreak\smallskip}
\pdfbookmark[0]{De nexu inter multitudinem classium. I}{nexus-prior}
\begin{center}
[I.]\par\medskip
\Large\bfseries DE NEXU INTER MULTITUDINEM CLASSIUM, IN QUAS FORMAE BINARIAE SECUNDI GRADUS DISTRIBUUNTUR, EARUMQUE DETERMINANTEM.\par\medskip
\normalsize COMMENTATIO PRIOR\par\medskip
SOCIETATI REGIAE EXHIBITA 1834\ .\ .\ .\ .
\end{center}
\BookRule{.18\textwidth}

% A544-P0279-D02
\Article{1}
Triginta tres iam elapsi sunt anni, ex quo principia nexus mirabilis, cui haec commentatio dicata est, deteximus, uti iam in fine \emph{Disquisitionum Arithmeticarum} annunciatum est. Sed aliae occupationes ab hac scrutatione per longum tempus detraxerant, donec recentiori tempore ad eam reverti et per novas curas eam ampliare contigit. Attamen quum haec nova Arithmeticae Sublimioris pars limites unius commentationis excedat, haecce prior formis determinantium negativorum dicata erit: formae vero determinantium positivorum, quae tractationem prorsus peculiarem requirunt, commentationi alteri reservatae manere debebunt.

% A544-P0279-D03
\Article{2}
Basis totius argumenti est disquisitio peculiaris circa multitudinem omnium combinationum valorum integrorum, quos duo numeri integri indefiniti $x,y$ intra ambitum praescriptum accipiunt. Manifesto hoc problema etiam sub aspectu geometrico exhiberi potest, ut eruatur multitudo \emph{numerorum complexorum}, quorum repraesentatio intra figuram praescriptam cadit. Indoles figurae ex indole lineae quae eam circumdat, adeoque pendebit vel ab unica aequatione inter coordinatas $x,y$ (quoties peripheria est curva in se rediens) vel a pluribus huiusmodi 
\SourceLeaf{280}%
% A544-P0280-D01
aequationibus (quoties .constat e pluribus partibus curvis seu rectis), pendebitque ab arbitrio nostro, utrum puncta numeris integris complexis respondentia, si quae forte in ipsa peripheria sint, multitudini annumerare velimus an inde excludere.

% A544-P0280-D02
In repraesentatione analytica problematis conditiones illius limitationis semper ita exhiberi poterunt, ut functio data variabilium $x,y$ vel una vel plures $P,Q,R$ etc. nancisci debeant valores positivos, vel non-negativos (prout valor $0$ vel excluditur vel admittitur).

% A544-P0280-D03
Ita e. g. si figura praescripta est circulus, cuius radius $=\surd A$, dum centrum cadit in punctum numero complexo integro respondens, conditio analytica erit, ut $A-xx-yy$ non sit negativus, siquidem, quod semper supponemus, puncta in ipsa peripheria sita retinere placet. Si figura est triangulum, tres functiones lineares $ax+by+c$, $a'x+b'y+c'$, $a''x+b''y+c''$ valores non-negativos habere debent, similiterque in aliis casibus.

% A544-P0280-D04
\Article{3}
Solutio problematis \emph{exacta}, generaliter loquendo, ita procedere debet, ut primo e natura conditionum variabilis altera e. g. $x$ intra limites coërceatur, inter quos valores singuli integri deinceps percurrant, et quot valores integri alterius $y$ singulis respondeant, eruere oportet, quorum multitudines dein in summam colligi debent. In casibus specialibus plerumque aderunt artificia specialia ad laborem abbreviandum.

% A544-P0280-D05
E. g. si figura, ut supra, est circulus, cuius radius $=\surd A$, sit $r$ integer proxime minor quam $\surd A$, vel ipse $\surd A$, si $A$ est quadratum. Perinde sint $r',r'',r'''$ etc. $r^{(r)}$ integri proxime minores quam $\surd(A-1)$, $\surd(A-4)$, $\surd(A-9)$ etc. usque ad $\surd(A-rr)$. Tunc multitudo quaesita erit
\[
\begin{aligned}
&=2r+1+2(2r'+1)+2(2r''+1)+2(2r'''+1)+\text{etc.}\\
&=1+4r+4r'+4r''+4r'''+\text{etc.}+4r^{(r)}
\end{aligned}
\]

% A544-P0280-D06
Brevior erit in hoc exemplo methodus sequens. Sit $q$ integer proxime minor quam $\surd\tfrac12 A$ (vel huic aequalis, quoties est integer), atque $r^{(q+1)}$, $r^{(q+2)}$, $r^{(q+3)}$ etc. integri proxime minores quam $\surd(A-(q+1)^2)$, $\surd(A-(q+2)^2)$, $\surd(A-(q+3)^2)$ etc. usque ad $\surd(A-rr)$. Tunc erit multitudo quaesita
\[
=4qq+1+4r+8(r^{(q+1)}+r^{(q+2)}+r^{(q+3)}+\text{etc.}+r^{(r)})
\]

